% -----------------------------------------------------------
\section{Priesthood}
% -----------------------------------------------------------
	\label{PRIESTHOOD}
	
% ----------------------
\subsection{See also}
% ----------------------

	\hyperref[LAW]{Law}, \hyperref[HUPOTASSO]{\grl{ὑποτάσσω}}, 

% % ----------------------
% \subsection{1828 American Dictionary of the English Language} % *1828
% % ----------------------
	% \label{PRIESTHOOD-1828}

	% \begin{compactitem}
		% \item
	% \end{compactitem}

% % ----------------------
% \subsection{Oxford English Dictionary} % *OED
% % ----------------------
	% \label{PRIESTHOOD-OED}

	% \begin{compactitem}
		% \item[]
	% \end{compactitem}

% ----------------------
\subsection{Scriptural Usage} % *SCRIPTURES
% ----------------------
	\label{PRIESTHOOD-SCRIPTURES}

	% % -----
	% \subsubsection{Old Testament} % *OT
	% % -----
		% \label{PRIESTHOOD-OT}
	
		% % --
		% \paragraph{KJV}
		% % --
	
			% \begin{tabular}{ll}
				% Total occurrences in the OT & 
			% \end{tabular}
			
			% \begin{compactdesc}
				% \item[]
			% \end{compactdesc}
			
		% % --
		% \paragraph{Hebrew Bible}
		% % --
			% \label{PRIESTHOOD-HEBREW}
		
			% %
			% \subparagraph{\texorpdfstring{\he{sample word}}{}}
			% %
				% \label{PAGA}
			
				% \begin{compactdesc}
					% \item[HALOT , PDF ] \textbf{}
						% \begin{compactitem}
							% \item[1] to let something hurt someone (Isaiah 53:6)
							% \item[2] to look after someone, intercede for, with \he{b:*} and \he{'El}/\he{`E:} (Isaiah 53:12)
							% \item[3] to urge someone
						% \end{compactitem}
					% \item[BDB , PDF ] \textbf{}
						% \begin{compactitem}
							% \item[1] cause to light upon (Isaiah 53:6) 
							% \item[2] cause one to entreat
							% \item[3] make entreaty
							% \item[4] interpose (Isaiah 53:12)
						% \end{compactitem}
					% \item[DCH PDF ] \textbf{}
						% \begin{compactitem}
							% \item[1] cause to come upon, visit (Isaiah 53:6)
							% \item[2] make entreaty (Isaiah 53:12)
							% \item[3] cause to entreat
						% \end{compactitem}
				% \end{compactdesc}
				
				% \begin{compactdesc}
					% \item[Some other OT uses] \textbf{}
						% \begin{compactdesc}
							% \item[]
						% \end{compactdesc}
				% \end{compactdesc}
			
		% % --
		% \paragraph{Septuagint}
		% % --
		
			% %
			% \subparagraph{\texorpdfstring{\gr{sample word}}{}}
			% %
		
	% -----
	\subsubsection{New Testament} % *NT
	% -----
		\label{PRIESTHOOD-NT}
	
		% --
		\paragraph{KJV}
		% --
	
			% \begin{tabular}{ll}
				% Total occurrences in the NT & 
			% \end{tabular}
			
			\begin{compactdesc}
				\item[{\hyperref[MATTHEW-1-17]{Matthew 1:17}}] \phantom{.}
					\begin{compactitem}
						\item See notes here also.
					\end{compactitem}
				\item[{\hyperref[MATTHEW-5-5]{Matthew 5:5}}] \phantom{.}
					\begin{compactitem}
						\item The language of ``inheriting'' makes me think of having a right to something by your connection to the original possessor, and that makes me think of the priesthood.
					\end{compactitem}
				\item[{\hyperref[MATTHEW-6-33]{Matthew 6:33}}] \phantom{.}
					\begin{compactitem}
						\item Note also other verses about the ``kingdom.''
					\end{compactitem}
				\item[{\hyperref[MATTHEW-7-24]{Matthew 7:24}}] \phantom{.}
					\begin{compactitem}
						\item See the footnote to ``rock'' especially.
					\end{compactitem}
				\item[{\hyperref[MATTHEW-8-8]{Matthew 8:8--10}}] \phantom{.}
					\begin{compactitem}
						\item See \hyperref[LUKE-7-6]{Luke 7:6--9} (including the use of \gr{τασσόμενος} in 8)
					\end{compactitem}
				\item[{\hyperref[MATTHEW-8-16]{Matthew 8:16}}]
				\item[{\hyperref[MATTHEW-8-24]{Matthew 8:24--27}}]
				\item[{\hyperref[MATTHEW-12-22]{Matthew 12:22--30}}] \phantom{.}
					\begin{compactitem}
						\item See also \hyperref[MARK-3-22]{Mark 3:22--27} and \hyperref[LUKE-11-14]{Luke 11:14--23}.
					\end{compactitem}
				\item[{\hyperref[MATTHEW-13-43]{Matthew 13:37--43}}] \phantom{.}
					\begin{compactitem}
						\item Notice the use of certain key words here, including ``good,'' \gr{δίκαιος}, \gr{ἀνομία}, and so on.
					\end{compactitem}
				\item[{\hyperref[MATTHEW-13-57]{Matthew 13:57}}] \phantom{.}
					\begin{compactitem}
						\item Notice the use of \gr{ἄτιμος}; perhaps connect to other verses that talk about ``honor'' in the sense of priesthood. It's a bit of a stretch, I think.
					\end{compactitem}
				\item[{\hyperref[MATTHEW-14-15]{Matthew 14:15--21}}] \phantom{.}
					\begin{compactitem}
						\item See this story as symbolizing the way that the priesthood and covenants are disseminated in a world: from the Father (Christ looks up into heaven), through Christ, to his disciples, and then to everyone else. Everyone benefits, but it happens in a particular order.
						\item Reminds me also of the new birth, where everyone is filled.
						\item This is worth exploring in a talk or CtR essay.
					\end{compactitem}
				\item[{\hyperref[MATTHEW-16-18]{Matthew 16:18--19}}]
				\item[{\hyperref[MATTHEW-17-14]{Matthew 17:14--21}}] \phantom{.}
					\begin{compactitem}
						\item I'm specifically wondering about the reason the disciples couldn't help the boy---I'm wondering if it's a lack of \gr{δύναμις}, which they might have remedied by ``prayer and fasting.'' This is a little speculative but it's interesting. See my notes at 17:21 and also note \hyperref[LUKE-9-43]{Luke 9:43}, \gr{ἐξεπλήσσοντο δὲ πάντες ἐπὶ τῇ μεγαλειότητι τοῦ θεοῦ}.
						\item Compare \hyperref[MARK-9-14]{Mark 9:14--29} and \hyperref[LUKE-9-38]{Luke 9:38--42}. Note the different ways that the ``prayer and fasting'' bit is different in the Greek.
					\end{compactitem}
				\item[{\hyperref[MATTHEW-21-18]{Matthew 21:18--22}}] \phantom{.}
					\begin{compactitem}
						\item Important for the relationship between priesthood and faith.
					\end{compactitem}
				\item[{\hyperref[MATTHEW-21-19]{Matthew 21:19--27}}] \phantom{.}
					\begin{compactitem}
						\item An important discussion of \gr{ἐξουσία}.
					\end{compactitem}
				\item[{\hyperref[MATTHEW-23-15]{Matthew 23:15}}] \phantom{.}
					\begin{compactitem}
						\item Kind of a weird connection but I was thinking of ``child of hell'' as like somewhat of a real relationship to the powers of hell.
					\end{compactitem}
				\item[{\hyperref[MARK-3-31]{Mark 3:31--35}}] \phantom{.}
					\begin{compactitem}
						\item The connection here is to the covenant/family relationship with Christ. He says something similar to this elsewhere too.
					\end{compactitem}
				\item[{\hyperref[MARK-5-24]{Mark 5:24--34}}] \phantom{.}
					\begin{compactitem}
						\item This is an important story because it shows that there is some sort of relationship between Christ's power and our faith. At least some of the time, that power is responsive to our faith. Christ says in verse 34, it's your faith that healed you, but that's not entirely true---it's his power that carried out the healing. But his power, at least some of the time, is responsive to our faith. He and his clothing were being touched by many people while he was walking through the crowd, but only one person approached him and touched him with the intent to be healed---only she exercised her faith in this way---and it was done. This is a really important story. In what ways are we ``touching'' Christ, rubbing shoulders with him and thronging him and even following him, without doing it with faith to make something happen?
						\item Compare the restriction on what Christ was able to do in \hyperref[MARK-6-5]{Mark 6:5--6}.
						\item See also \hyperref[MARK-6-55]{Mark 6:55-56}.
					\end{compactitem}
				\item[{\hyperref[MARK-7-32]{Mark 7:32--35}}] \phantom{.}
					\begin{compactitem}
						\item Another example where a healing involved a lot of physical touch. But note in the story before this one, Christ casts out a spirit from afar.
					\end{compactitem}
				\item[{\hyperref[MARK-14-62]{Mark 14:62}}]
				\item[{\hyperref[LUKE-4-36]{Luke 4:36}}] \phantom{.}
					\begin{compactitem}
						\item See also other verses in event \#30 of Christ's life.
					\end{compactitem}
				\item[{\hyperref[LUKE-5-17]{Luke 5:17--26}}]
					\begin{compactitem}
						\item I think the key question is in verse 23---which of these is actually easier? They require different powers, and note that we have gained enough knowledge to make some of the healing power available to ourselves now.
					\end{compactitem}
				\item[{\hyperref[LUKE-7-2]{Luke 7:2--10}}]
				\item[{\hyperref[LUKE-8-4]{Luke 8:4--18}}] \phantom{.}
					\begin{compactitem}
						\item This is Luke's version of the parable of the sower. The fact that the seed is the \gr{λόγος} and gets planted in us (compare \hyperref[ALMA-32]{Alma 32}) could be relevant to the priesthood, since Christ as \gr{λόγος}.
					\end{compactitem}
				\item[{\hyperref[LUKE-10-17]{Luke 10:17--20}}] \phantom{.}
					\begin{compactitem}
						\item Note verse 19 in particular, where Christ gives them \gr{ἐξουσία} over the \gr{δύναμις} of the ``enemy.'' Important passage.
					\end{compactitem}
				\item[{\hyperref[LUKE-10-22]{Luke 10:22}}] \phantom{.}
				\item[{\hyperref[LUKE-11-33]{Luke 11:33--36}}] \phantom{.}
					\begin{compactitem}
						\item See my footnotes on these verses, and compare to \hyperref[DC88-5]{DC 88:5ff}.
					\end{compactitem}
				\item[{\hyperref[LUKE-12-42]{Luke 12:42--44}}] \phantom{.}
				\item[{\hyperref[LUKE-17-20]{Luke 17:20--21}}] \phantom{.}
				\item[{\hyperref[LUKE-20-14]{Luke 20:14}}] \phantom{.}
					\begin{compactitem}
						\item Plus other aspects of this parable.
					\end{compactitem} 
				\item[{\hyperref[LUKE-22-53]{Luke 22:53}}] \phantom{.}
				\item[{\hyperref[LUKE-22-69]{Luke 22:69}}] \phantom{.}
				\item[{\hyperref[LUKE-24-49]{Luke 24:49}}] \phantom{.}
				\item[{\hyperref[JOHN-4-7]{John 4:7--14}}] \phantom{.}
				\item[{\hyperref[JOHN-5-17]{John 5:17--47}}] \phantom{.}
					\begin{compactitem}
						\item This is a long block of verses, and not all of them are equally important; some, like \hyperref[JOHN-5-26]{verse 26} for example, are extremely important.
					\end{compactitem}
				\item[{\hyperref[JOHN-6-32]{John 6:32--35}}] \phantom{.}
					\begin{compactitem}
						\item More verses besides these are relevant here--the entire bread of life discussion is important, and especially verses 56--57. But the entire discussion deserves some study in relation to the priesthood.
					\end{compactitem}
				\item[{\hyperref[JOHN-10-17]{John 10:17--18}}] \phantom{.}
				\item[{\hyperref[JOHN-14-6]{John 14:6--11}}] \phantom{.}
				\item[{\hyperref[JOHN-14-16]{John 14:16--20}}] \phantom{.}
				\item[{\hyperref[JOHN-15-1]{John 15:1--11}}] \phantom{.}
					\begin{compactitem}
						\item Some of this just seems to be so clearly talking about the covenant relationship we enter into with Christ, and the structure of the priesthood relationship we create with him through covenants. Look at 15:4, for example, or verse 10.
						\item Now that I'm reading the rest of the chapter, I think other parts are relevant, too, likes verses 15--16.
							\begin{compactitem}
								\item Though note that ``ordained'' in 16 is \gr{ἔθηκα}, a form of \gr{τίθημι}, so you'd have to check this out carefully.
								\item But 16 clearly relates to what Christ says earlier in the chapter about being part of him and having that relationship.
							\end{compactitem}
						\item It's almost like, if there is anything like an oath and covenant of the priesthood in the NT, it looks like this chapter.
					\end{compactitem}
				\item[{\hyperref[JOHN-16-15]{John 16:15}}] \phantom{.}
					\begin{compactitem}
						\item Goes with the previous chapter too.
					\end{compactitem}
				\item[{\hyperref[JOHN-17-2]{John 17:2}}] \phantom{.}
					\begin{compactitem}
						\item There are other verses in this chapter that are also relevant to priesthood, including a few things about the indwelling relationship.
					\end{compactitem}
				\item[{\hyperref[JOHN-19-10]{John 19:10--11}}]
				\item[{\hyperref[JOHN-20-21]{John 20:21--22}}]
				\item[{\hyperref[ROMANS-1-16]{Romans 1:16--20}}]
				\item[{\hyperref[ROMANS-8-7]{Romans 8:7}}]
				\item[{\hyperref[ROMANS-13-1]{Romans 13:1--5}}] \phantom{.}
					\begin{compactitem}
						\item Very interesting verses, especially the first three. Possible connections to DC 88 as well as to the structure of the priesthood, with God at the top, and everything lined up below him. The word \gr{τάσσω} gets used in these verses and that word is important---notice also that the KJV translation of verse 2 is misleading. The second and third ``resist''s are \gr{ἀνθίστημι}, but the first is a participle of \gr{ἀντιτάσσω}. That word does mean ``resist'' but Paul is clearly picking up on the use of \gr{τάσσω} in the first verse. Then ``ordinance'' is \gr{διαταγή}, which obviously comes from \gr{διατάσσω}, picking up again on that idea. 
						\item For studying the priesthood more generally, it would be helpful to pick up the importance of \gr{τάσσω} in various places, to see if we could make a connection. Actually, now that I'm looking, I think \gr{τάσσω} only appears around eight times in the NT, but the use I was thinking of is important---it's \hyperref[LUKE-7-8]{Luke 7:8}.
					\end{compactitem}
				\item[{\hyperref[2CORINTHIANS-3-6]{2 Corinthians 3:6--11}}]
				\item[{\hyperref[2CORINTHIANS-12-9]{2 Corinthians 12:9--10}}]
				\item[{\hyperref[2CORINTHIANS-13-4]{2 Corinthians 13:4}}] \phantom{.}
					\begin{compactitem}
						\item There are other verses around here that talk about other things connected to \gr{δύναμις}, in various forms of the word family.
					\end{compactitem}
				\item[{\hyperref[2CORINTHIANS-13-10]{2 Corinthians 13:10}}] \phantom{.}
					\begin{compactitem}
						\item And here you get a reference to \gr{ἐξουσία}.
					\end{compactitem}
				\item[{\hyperref[GALATIANS-3-19]{Galatians 3:19}}]
				\item[{\hyperref[EPHESIANS-1-17]{Ephesians 1:17--23}}]
				\item[{\hyperref[EPHESIANS-2-18]{Ephesians 2:18}}]
				\item[{\hyperref[EPHESIANS-3]{Ephesians 3}}] \phantom{.}
					\begin{compactitem}
						\item Not every verse, but \gr{δύναμις} appears and so does \gr{ἐξουσία}, and it talks about God's ``workings.'' I'd have to go through it more carefully to pull out the specific verses but there is some relevant stuff here.
					\end{compactitem}
				\item[{\hyperref[EPHESIANS-4-6]{Ephesians 4:6--16}}] \phantom{.}
					\begin{compactitem}
						\item This is a very powerful set of verses when read carefully, in the light of the law of Christ and the priesthood.
					\end{compactitem}
				\item[{\hyperref[PHILIPPIANS-2-5]{Philippians 2:5--11}}] \phantom{.}
					\begin{compactitem}
						\item I'm not sure what to make of these verses. 
						\item The interpretation of these verses among scholars is controversial and very little is settled.
					\end{compactitem}
				\item[{\hyperref[PHILIPPIANS-3-8]{Philippians 3:8--14}}] \phantom{.}
					\begin{compactitem}
						\item Mentions the power of the resurrection.
					\end{compactitem}
				\item[{\hyperref[COLOSSIANS-1-9]{Colossians 1:9--29}}] \phantom{.}
					\begin{compactitem}
						\item Especially verse 13, which mentions \gr{ἐκ τῆς ἐξουσίας τοῦ σκότους}.
						\item Also verse 17; there are a lot of verses relevant to the priesthood here.
					\end{compactitem}
			\end{compactdesc}
			
		% % --
		% \paragraph{Greek New Testament}
		% % --
			% \label{PRIESTHOOD-GREEK}
		
			% %
			% \subparagraph{\texorpdfstring{\gr{sample word}}{}}
			% %
				% \label{KATALLAGE}
			
				% \begin{compactdesc}
					% \item[BDAG , PDF ] \textbf{}
						% \begin{compactitem}
							% \item 
						% \end{compactitem}
				% \end{compactdesc}
				
				% \begin{compactdesc}
					% \item[Some other NT uses] \textbf{}
						% \begin{compactdesc}
							% \item[]
						% \end{compactdesc}
				% \end{compactdesc}
		
	% -----
	\subsubsection{Book of Mormon} % *BOM
	% -----
		\label{PRIESTHOOD-BOM}
	
		% \begin{tabular}{ll}
			% Total occurrences in the BOM & 
		% \end{tabular}
		
		\begin{compactdesc}
			\item[{\hyperref[1NEPHI-2-14]{1 Nephi 2:14}}]
			\item[{\hyperref[1NEPHI-11-31]{1 Nephi 11:31}}]	
			\item[{\hyperref[1NEPHI-14-14]{1 Nephi 14:14}}]	
			\item[{\hyperref[1NEPHI-15-15]{1 Nephi 15:15}}] \phantom{.}
				\begin{compactitem}
					\item There's a way of understanding the ``true vine'' stuff that connects you to the parts in John about the vine. I think you can understand all of that discussion as being about the way that power comes from God to Christ to us.
				\end{compactitem}
			\item[{\hyperref[1NEPHI-17-47]{1 Nephi 17:47--54}}]
			\item[{\hyperref[2NEPHI-6-2]{2 Nephi 6:2}}]
			\item[{\hyperref[2NEPHI-19-6]{2 Nephi 19:6--7}}]
			\item[{\hyperref[JACOB-4-6]{Jacob 4:6--9}}] \phantom{.}
				\begin{compactitem}
					\item Not really verse 8, but it picks up again in verse 9.
				\end{compactitem}
			\item[{\hyperref[MOSIAH-3-17]{Mosiah 3:17}}] \phantom{.}
				\begin{compactitem}
					\item I want to read this verse in light of \gr{λόγος}-type stuff, where there's some fundamental connection between Christ and reality, and that's the reason why there will be no other way nor means for achieving salvation.
				\end{compactitem}
			\item[{\hyperref[MOSIAH-15-1]{Mosiah 15:1--5}}]
			\item[{\hyperref[MOSIAH-15-1]{Mosiah 15:11}}] \phantom{.}
				\begin{compactitem}
					\item And possibly other verses here and in chapter 14 about Christ's ``seed.'' To be his seed is to be the heir of the kingdom of God? We need that lineage through the priesthod?
				\end{compactitem}
			\item[{\hyperref[MOSIAH-23-17]{Mosiah 23:17}}] \phantom{.}
				\begin{compactitem}
					\item I'm sure there are other uses of ``consecrate'' that I haven't picked up on in this list.
				\end{compactitem}
			\item[{\hyperref[MOSIAH-29-19]{Mosiah 29:19--20}}] \phantom{.}
			\item[{\hyperref[ALMA-4-20]{Alma 4:20}}]
			\item[{\hyperref[ALMA-5-49]{Alma 5:49}}] \phantom{.}
				\begin{compactitem}
					\item See also verse 44 here, and 54.
					\item 54 is particularly important because it notes that the people who have been baptized have been received into the order. See also Alma 7:22 below, and \hyperref[MOSES-6-66]{Moses 6:66--68}.
				\end{compactitem}
			\item[{\hyperref[ALMA-6-1]{Alma 6:1}}] \phantom{.}
				\begin{compactitem}
					\item See also verse 8.
				\end{compactitem}
			\item[{\hyperref[ALMA-7-22]{Alma 7:22}}] \phantom{.}
				\begin{compactitem}
					\item This verse is very important, because it makes clear that when we're following Christ and have made covenants, we are part of the ``holy order of God,'' or the priesthood.
					\item And this is true for every person, no matter who you are.
				\end{compactitem}
			\item[{\hyperref[ALMA-8-4]{Alma 8:4}}]
			\item[{\hyperref[ALMA-8-24]{Alma 8:24}}]
			\item[{\hyperref[ALMA-8-31]{Alma 8:31--32}}] 
			\item[{\hyperref[ALMA-13]{Alma 13}}] \phantom{.}
				\begin{compactitem}
					\item Maybe not the entire chapter, but a lot of it.
					\item Note that, in accordance with what happens at the end of Alma 12 (see my notes on the \hyperref[EPIST-REVELATION-BOM]{epistemology of revelation}), the priesthood is closely connected to the way God disseminates knowledge of the plan of salvation to people.
					\item See Helaman 8:16ff below.
				\end{compactitem}
			\item[{\hyperref[ALMA-14-25]{Alma 14:25--29}}] \phantom{.}
				\begin{compactitem}
					\item Notice the connection to faith specifically.
				\end{compactitem}
			\item[{\hyperref[ALMA-29-13]{Alma 29:13}}]
			\item[{\hyperref[ALMA-39-15]{Alma 39:15--19}}]
			\item[{\hyperref[ALMA-43-2]{Alma 43:2}}]
			\item[{\hyperref[HELAMAN-5-11]{Helaman 5:11--12}}]
			\item[{\hyperref[HELAMAN-5-18]{Helaman 5:18}}]
			\item[{\hyperref[HELAMAN-5-23]{Helaman 5:23}}] \phantom{.}
				\begin{compactitem}
					\item This is the start of Nephi and Lehi's experience in the prison.
				\end{compactitem}
			\item[{\hyperref[HELAMAN-8-16]{Helaman 8:16}}]
			\item[{\hyperref[HELAMAN-10-1]{Helaman 10:1--11}}]
			\item[{\hyperref[HELAMAN-10-15]{Helaman 10:15--17}}]
			\item[{\hyperref[3NEPHI-5-13]{3 Nephi 5:13}}]
			\item[{\hyperref[3NEPHI-8-1]{3 Nephi 8:1}}]
			\item[{\hyperref[3NEPHI-9-15]{3 Nephi 9:15}}]
			\item[{\hyperref[3NEPHI-11-21]{3 Nephi 11:21--28}}]
			\item[{\hyperref[3NEPHI-12-1]{3 Nephi 12:1}}]
			\item[{\hyperref[3NEPHI-15-9]{3 Nephi 15:9--10}}]
			\item[{\hyperref[3NEPHI-18-36]{3 Nephi 18:36--38}}]
sv			\item[{\hyperref[3NEPHI-18-5]{3 Nephi 18:5}}]
			\item[{\hyperref[3NEPHI-20-22]{3 Nephi 20:22}}]
			\item[{\hyperref[3NEPHI-21-25]{3 Nephi 21:25}}]
			\item[{\hyperref[3NEPHI-28-10]{3 Nephi 28:20}}]
			\item[{\hyperref[3NEPHI-28-39]{3 Nephi 28:39}}]
			\item[{\hyperref[MORMON-8-24]{Mormon 8:24}}]
			\item[{\hyperref[MORMON-9-13]{Mormon 9:13}}]
			\item[{\hyperref[MORMON-9-17]{Mormon 9:17}}]
			\item[{\hyperref[ETHER-3-14]{Ether 3:14}}]
			\item[{\hyperref[ETHER-12-10]{Ether 12:10}}]
			\item[{\hyperref[ETHER-12-28]{Ether 12:28}}] \phantom{.}
				\begin{compactitem}
					\item I'm reading the ``fountain of all righteousness'' it in a priesthood way.
				\end{compactitem}
			\item[{\hyperref[MORONI-2-2]{Moroni 2:2--3}}]
			\item[{\hyperref[MORONI-3-2]{Moroni 3:2--4}}]
			\item[{\hyperref[MORONI-7-27]{Moroni 7:27}}]
		\end{compactdesc}
		
	% -----
	\subsubsection{Doctrine and Covenants} % *DC
	% -----
		\label{PRIESTHOOD-DC}
	
		% \begin{tabular}{ll}
			% Total occurrences in the DC & 
		% \end{tabular}
		
		\begin{compactdesc}
			\item[{\hyperref[DC18-23]{DC 18:23--25}}]
			\item[{\hyperref[DC19-1]{DC 19:1--3}}] \phantom{.}
				\begin{compactitem}
					\item Especially the ``subdue'' language; see the footnotes.
				\end{compactitem}
			\item[{\hyperref[DC28-7]{DC 28:7}}] \phantom{.}
				\begin{compactitem}
					\item This verse belongs in a group that also includes \hyperref[DC35-18]{DC 35:18}, \hyperref[DC64-5]{DC 64:5}, and \hyperref[DC84:19]{DC 84:19}, where the revelations are talking about the ``keys'' and ``mysteries'' together.
				\end{compactitem}
			\item[{\hyperref[DC29-29]{DC 29:29--41}}]
			\item[{\hyperref[DC41-3]{DC 41:3--5}}]
			\item[{\hyperref[DC45-8]{DC 45:8}}] 
			\item[{\hyperref[DC49-5]{DC 49:5--6}}] 
			\item[{\hyperref[DC50-27]{DC 50:27--28}}]
			\item[{\hyperref[DC53-3]{DC 53:3}}]
			\item[{\hyperref[DC61-1]{DC 61:1}}]
			\item[{\hyperref[DC63-59]{DC 63:59}}]
			\item[{\hyperref[DC66-2]{DC 66:2}}]
			\item[{\hyperref[DC68-14]{DC 68:14--21}}] \phantom{.}
				\begin{compactitem}
					\item Note that not all of these verses were in the original version of section 68; a lot of it was added later.
					\item Verse 17 is really important though, as I think it comes to shape other ways of seeing the priesthood later too.
					\item Verse 20 is important too.
				\end{compactitem}
			\item[{\hyperref[DC75]{DC 75}}] \phantom{.}
				\begin{compactitem}
					\item See especially the \hyperref[DC75-HB]{historical background}, which talks about the ``high priesthood.''
				\end{compactitem} 
			\item[{\hyperref[DC76]{DC 76}}] \phantom{.}
				\begin{compactitem}
					\item There are a number of verses in this section, like 42, which describe the relationship between the Father and the Son, and I think those verses are important for getting at the structure of the priesthood and how those figures relate to it.
					\item See also \hyperref[DC76-50]{DC 76:50--59}, which are very important.
					\item See \hyperref[DC76-86]{86--88} and notes there.
					\item See 95, and 107--108.
				\end{compactitem}
			\item[{\hyperref[DC77-3]{DC 77:3}}]
			\item[{\hyperref[DC77-11]{DC 77:11}}]
			\item[{\hyperref[DC78]{DC 78}}] \phantom{.}
				\begin{compactitem}
					\item \hyperref[DC78-1]{DC 78:1--2}
					\item \hyperref[DC78-8]{DC 78:8}
					\item \hyperref[DC78-16]{DC 78:16}
					\item This section has some subtle but interesting references to ``order'' and other ideas connected to the priesthood.
				\end{compactitem}
			\item[{\hyperref[DC81-1]{DC 81:1--2}}]
			\item[{\hyperref[DC82-3]{DC 82:3--4}}]
			\item[{\hyperref[DC82-10]{DC 82:10, 15}}]
			\item[{\hyperref[DC82-17]{DC 82:17--21}}]
			\item[{\hyperref[DC84-1]{DC 84:1--61}}] \phantom{.}
				\begin{compactitem}
					\item This section is full of very important things about the priesthood, including the ``sonship'' relationship used to describe a certain relation in the priesthood.
					\item Verse 18 calls the Melchizedek priesthood ``the priesthood which is after the holiest order of God.''
					\item See also \hyperref[DC84-63]{verse 63}.
				\end{compactitem}
			\item[{\hyperref[DC85-11]{DC 85:11}}]
			\item[{\hyperref[DC86]{DC 86}}] \phantom{.}
				\begin{compactitem}
					\item This whole section is important for the priesthood; see \hyperref[DC86-1]{verse 1} and my notes there.
					\item This section has two halves---the commentary on the wheat and the tares, and the comments about priesthood. The connection between them is not super clear.
					\item I read the second half as Christ speaking to his original apostles.
				\end{compactitem}
			\item[{\hyperref[DC88-5]{DC 88:5--13}}] \phantom{.}
				\begin{compactitem}
					\item These verses are fundamental for understanding the relationship between God and Christ, and the priesthood. See also \hyperref[DC107-3]{DC 107:3}, which is saying basically the same thing as verse 5 here.
					\item To me it seems clear that the ``light'' being spoken of in verses 12--13, and possibly the others verses, is the priesthood. It is the ``law by which all things are governed, even the power of God.'' This gives you a much more direct connection to BY-style ideas of the priesthood than you can find elsewhere.
				\end{compactitem} 
			\item[{\hyperref[DC88-14]{DC 88:14--60?}}] \phantom{.} 
				\begin{compactitem}
					\item All the verses that discuss ``laws'' here in this section are relevant to the priesthood, and are very important.
				\end{compactitem}
			\item[{\hyperref[DC93]{DC 93}}] \phantom{.} 
				\begin{compactitem}
					\item Another whole section that is relevant to many different aspects of the priesthood.
					\item \hyperref[DC93-20]{Verse 20} is particularly good, in thinking about \hyperref[DC107-3]{DC 107:3} and why the priesthood has that name. See also \hyperref[DC93-22]{DC 93:22} and \hyperref[MOSIAH-5-1]{Mosiah 5:1--10}.
				\end{compactitem}
			\item[{\hyperref[DC94-6]{DC 94:6}}]
			\item[{\hyperref[DC98-3]{DC 98:3}}]
			\item[{\hyperref[DC98-14]{DC 98:14--15}}]
			\item[{\hyperref[DC101-77]{DC 101:77--78}}]
			\item[{\hyperref[DC106]{DC 106}}] \phantom{.}
				\begin{compactitem}
					\item This section has a few different instances of royal-related language, including ``crown,'' ``kingdom,'' and ``scepter''
				\end{compactitem}
			\item[{\hyperref[DC107-1]{DC 107:1--20}}]
			\item[{\hyperref[DC107-40]{DC 107:40--57}}]
			\item[{\hyperref[DC110-11]{DC 110:11--16}}] \phantom{.}
				\begin{compactitem}
					\item This is the second part of this section, which records visitations and various ``keys'' being given.
				\end{compactitem}
			\item[{\hyperref[DC112-28]{DC 112:28--33}}]
			\item[{\hyperref[DC113-6]{DC 113:6}}]
			\item[{\hyperref[DC113-8]{DC 113:8}}]
			\item[{\hyperref[DC115-19]{DC 115:19}}]
			\item[{\hyperref[DC121-4]{DC 121:4}}]
			\item[{\hyperref[DC121-21]{DC 121:21}}]
			\item[{\hyperref[DC121-29]{DC 121:29--32}}]
			\item[{\hyperref[DC121-34]{DC 121:34--46}}]
			\item[{\hyperref[DC124-19]{DC 124:19}}]
			\item[{\hyperref[DC124-28]{DC 124:28}}]
				\begin{compactitem}
					\item This verse talks about the ``fulness of the priesthood.''
				\end{compactitem}
			\item[{\hyperref[DC124-33]{DC 124:33--34}}]
			\item[{\hyperref[DC124-42]{DC 124:42}}]
			\item[{\hyperref[DC124-55]{DC 124:55}}]
			\item[{\hyperref[DC124-91]{DC 124:91--97}}] \phantom{.}
				\begin{compactitem}
					\item Talking about the ``patriarchal'' priesthood here?
				\end{compactitem}
			\item[{\hyperref[DC124-123]{DC 124:123--145}}] \phantom{.}
				\begin{compactitem}
					\item A lot of great stuff here, including very interesting references to the ``holy spirit of promise'' in 124.
					\item You have to read the first few verses in this block in light of what is said about Hyrum earlier, in verses 91--97.
				\end{compactitem}
			\item[{\hyperref[DC127-6]{DC 127:6--8}}] \phantom{.}
				\begin{compactitem}
					\item Especially verse 8; remember, this is written in September 1842, after the restoration of the endowment ceremony, but before the other ordinances in 1843.
				\end{compactitem}
			\item[{\hyperref[DC128-8]{DC 128:8--18}}]
			\item[{\hyperref[DC128-20]{DC 128:20--21}}]
			\item[{\hyperref[DC128-23]{DC 128:23}}]
			\item[{\hyperref[DC130-20]{DC 130:20--21}}]
			\item[{\hyperref[DC132]{DC 132}}] \phantom{.}
				\begin{compactitem}
					\item Lots of stuff in this section relevant to the priesthood and laws; too much to discuss here, at least for now.
					\item Verse 28 is important; reminds me of the opening of \hyperref[DC107]{DC 107}.
				\end{compactitem}
			\item[{\hyperref[DC138-51]{DC 138:51}}]
		\end{compactdesc}
		
	% % -----
	% \subsubsection{Pearl of Great Price} % *PGP
	% % -----
		% \label{PRIESTHOOD-PGP}
	
		% \begin{tabular}{ll}
			% Total occurrences in the PGP & 
		% \end{tabular}
		
		% \begin{compactdesc}
			% \item[]
		% \end{compactdesc}
		
% ----------------------
\subsection{Key Phrases} % *KEYPHRASES *PHRASES
% ----------------------
	\label{PRIESTHOOD-PHRASES}

	\begin{compactdesc}
		
		\item[fulness of (the) priesthood] \textbf{1} | \hyperref[DC124-28]{DC 124:28}
			\begin{compactdesc}
				\item[Bible anchor verses] ---
					% \begin{compactdesc}
						% \item[Romans 5:5] \gr{}
					% \end{compactdesc}
				\item[Commentary] This phrase occurs elsewhere in Joseph Smith's sermons, especially in the context of the second anointing and further temple ordinances in 1843.
			\end{compactdesc}
		
		% \item[] \textbf{\#times} | list
			% % \begin{compactdesc}
				% % \item[Bible anchor verses] \phantom{.}
					% % \begin{compactdesc}
						% % \item[Romans 5:5] \gr{}
					% % \end{compactdesc}
				% % \item[Commentary] ---
			% % \end{compactdesc}
			
		% \item[] \textbf{\#times} | list
			% % \begin{compactdesc}
				% % \item[Bible anchor verses] \phantom{.}
					% % \begin{compactdesc}
						% % \item[Romans 5:5] \gr{}
					% % \end{compactdesc}
				% % \item[Commentary] ---
			% % \end{compactdesc}
			
	\end{compactdesc}
	
% % ----------------------
% \subsection{Bible Dictionaries} % *DICTIONARIES
% % ----------------------
	% \label{PRIESTHOOD-BD}

% ----------------------
\subsection{Restoration Usage} % *RESTORATION
% ----------------------
	\label{PRIESTHOOD-RESTORATION}

	% -----
	\subsubsection{Joseph Smith} % *JS
	% -----
		\label{PRIESTHOOD-JS}
	
		\begin{compactdesc}
			\item[{\hyperref[EMS-1832-JS-CIR-SUM-1832]{Circa Summer 1832}}] A History of the life of Joseph Smith J\uline{r}. an account of his marvilous experience and of all the mighty acts which he doeth in the name of Jesus Ch[r]ist the son of the living God of whom he beareth record and also an account of the rise of the church of Christ in the eve of time according as the Lord brought forth and established by his hand <firstly> he receiving the testamony from on high seccondly the ministering of Angels thirdly the reception of the holy Priesthood by the ministring of---Aangels to adminster the letter of the \sout{Law} <Gospel---> <---the Law and commandments as they were given unto him---> and \sout{in} <the> ordinencs, forthly a confirmation and reception of the high Priesthood after the holy order of the son of the living God power and ordinence from on high to preach the Gospel in the administration and demonstration of the spirit the Kees of the Kingdom of God confered upon him and the continuation of the blessings of God to him \&c------
			\item[{\hyperref[EMS-1837-JS-6-APR-1837]{6 April 1837}}] Joseph Smith jr. rose and spoke on the subject of the Priesthood. \hl{The Melchisedec High priesthood, he said was no other then the priesthood of the Son of God. There are certain ordinances which belong to the priesthood, and certain results flow from it.} The presidents, or presidency are over the church, and revelations of the mind and will of God to the church are to come through the presidency. This is the order of heaven and the power and privilege of this priesthood. It is also the privilege of any officer in this church, to obtain revelations so far as relates to his particular calling or duty in the church. All are bound by the principles of virtue and happiness, but one great privilege of this priesthood is to obtain revelations, as before observed, of the mind and will of God. It is also the privilege of the Melchisedec priesthood, to reprove, rebuke and admonish, as well as to receive revelations.
			\item[{\hyperref[EMS-1841-JS-5-JAN-1841-WC]{5 January 1841}}] \hl{All priesthood is Melchizedeck; but there are different portions or degrees of it.} That portion which brought Moses to speak with God face to face was taken away; but that which brought the ministry of angels remained. All the Prophets had the Melchizedeck Priesthood and was ordained by God himself.
			\item[{\hyperref[EMS-1842-JS-15-JUL-1842-A]{15 July 1842}}] The Lord has at various times commenced this kind of government, and tendered his services to the human family. He selected Enoch, whom he directed, and gave his law unto, and to the people who were with him; and when the world in general would not obey the commands of God, after walking with God, he translated Enoch and his church, \hl{and the priesthood or government of heaven}, was taken away.
			\item[{\hyperref[EMS-1842-JS-27-JUL-1842]{27 July 1842}}] \phantom{.}
				\begin{compactitem}
					\item From a revelation regarding plural marriage; many almost off-hand but important comments about how they viewed the priesthood.
				\end{compactitem}
			\item[{\hyperref[EMS-1843-JS-23-MAR-1843-SAW]{23 March 1843}}] \textbf{these promises I Seal upon all of their heads in the name of Jesus Christ \hl{by the Law of the holy preisthood} Even so Amen}
		\end{compactdesc}
	
	% -----
	\subsubsection{Brigham Young} % *BY
	% -----
		\label{PRIESTHOOD-BY}
	
		\begin{compactdesc}
			\item[{\hyperref[EMS-1854-BY-12-FEB-1854]{12 February 1854}}] I recolect once when I was preaching a question was asked me, ``What is \sout{the} Preisthood.'' The answer was ready, and perfectly simple in its nature, \sout{and} plain to be understood, <and> \sout{\&} \sout{is} couched in a a short sentence (viz) \hl{Preisthood is a perfect system of Government that rules and reigns in eternity.} A question at once arises in the mind; ``\uline{where} \uline{is} \uline{eternity}.'' The answer is at hand, \uline{Eternity} is here;---we are in eternity just as much so as any other \sout{created} beings in heaven or on earth, Heavenly beings are no more in eternity than we are. We are in the midst of eternity; and \hl{when we become aquainted with the system of Gov}\supers{\hl{t.}}\hl{ and the laws that rule in eternity, we shall then know it callculated to endure; to govern and control all things which are in heaven, and on the earth. The eternal preisthood of God,---the Gov}\supers{\hl{t.}}\hl{ of God,---the laws of eternity, <is> }\sout{are} \hl{a pure and perfect system of Gov}\supers{\hl{t.}}
			\item[{\hyperref[EMS-1854-BY-27-JUN-1854]{27 June 1854}}] \hl{The Priesthood is a pure system of government which governs the world and a small portion has come down to earth. It is the constitutional laws and pure system by which men are to be governed.}
			\item[{\hyperref[EMS-1854-BY-19-NOV-1854]{19 November 1854}}] There is no kingdom without a space, and there is no space without a kingdom says the prophet Joseph. \hl{Now let me tell you there is no world made without that world has motion, and it is governed and controled by the law of the eternal preisthood; and if you want to know what preisthood is, it is a perfect system of Government by which planets and all creations are governed.}
			\item[{\hyperref[EMS-1854-BY-3-DEC-1854]{3 December 1854}}] There is a law that governs man thus far; but the law of the celestial kingdom, as I have frequently told you, is, and always will be, the same to all the children of Adam. \hl{When we talk of the celestial law which is revealed from heaven, that is, the Priesthood, we are talking about the principle of salvation, a perfect system of government, of laws and ordinances, by which we can be prepared to pass from one gate to another, and from one sentinel to another, until we go into the presence of our Father and God.} This law has not always been upon the earth; and in its absence, other laws have been given to the children of men for their improvement, for their education, for their government, and to prove what they would do when left to control themselves; and what we now call tradition has grown out of these circumstances.
			\item[{\hyperref[EMS-1859-BY-31-JUL-1859]{31 July 1859}}] \hl{Eternal existence depends solely upon adopting and carrying out in our lives the principles couched in the term ``holy Priesthood,'' which alone tend to life and eternal duration and exaltation...The holy Priesthood is a system of laws and government that is pure and holy; and if it is adhered to by intelligent man, whom God has created a little lower than angels, it is calculated to preserve our tabernacles in eternal being; otherwise they will be resolved into native element...Nothing can satisfy, except being perfectly subject to the law that will preserve them in their identity to all eternity, and that is the holy Priesthood...It is by the law of the Priesthood that men are, and by that law they may maintain their eternal identity.} A strict observance of those laws will secure an inheritance in that kingdom where death never enters, and all else will sooner or later pass away as a night vision...\hl{We have to learn and practice eternal principles, to obtain eternal life; and they are the principles of the holy Priesthood.} God has given man an agency, and it behooves us to understand and practice the principles of life---to live our religion and walk humbly with our God, living according to the laws and regulations of the holy Priesthood so far as it is revealed. The principles of eternal life that are set before us are calculated to exalt us to power and preserve us from decay. If we choose to take the opposite course and to imbibe and practice the principles that tend to death, the fault is with ourselves...\hl{Seek for the principles that pertain to eternal life---the principles of the holy Priesthood.}
			\item[{\hyperref[EMS-1864-BY-31-JUL-1864]{31 July 1864}}] \hl{...we frequently hear people inquire what the priesthood is it is a pure holy system of government the priesthood is the law that governs and controls the heavenly bodies the priesthood is the power and the principle that holds all things in existence by the word of him who is the author of all things}...this priesthood which God has revealed is for the salvation of the children of man. \hl{This priesthood which is a perfect system of government} is calculated to govern and control eventually the earth the inhabitants of the earth and all things pertaining to the earth...
			\item[{\hyperref[EMS-1865-BY-6-JUN-1865]{6 June 1865}}] The priesthood is the power of God; it rules in the eternal world, and should be recognized by all the inhabitants of the earth.
			\item[{\hyperref[EMS-1866-BY-17-JUN-1866]{17 June 1866}}] ...the kingdom of God always takes care of itself] it is a living moving thing [space] \hl{the priesthood of the Son of God comprises in its operations the kingdom of God if we were to ask what the priesthood of the Son of God is I know of no better language used what has been conveyed to me by the Spirit that is to say to the people it is a pure system of government}] those let[?] those that are governed live according to this system of laws of this pure government don't you see that all harmonizes and each and every part of this kingdom or the members of the kingdom remains one with the other...it is not any man that sits[?] upon this earth that carries on this work it guides itself it directs itself by the great machinist whose work and workmanship we are whose kingdom this is it is all governed and controlled by him who knows all things] and he will bring forth the and he will bring forth the righteous and the just [inserted above line: and] humble [inserted above line: the meek of the earth all] those that love to serve him and keep his commandments and enjoy the fullness of his glory [space] these ordinances this kingdom and work is proffered to the whole human family all who desire may accept all the kingdom upon the terms of strict obedience...
			\item[{\hyperref[EMS-1870-BY-30-OCT-1870]{30 October 1870}}] We are called, as it has been told you, to redeem the nations of the earth. The fathers cannot be made perfect without us; we cannot be made perfect without the fathers. There must be this chain in the holy Priesthood; it must be welded together from the latest generation that lives on the earth back to Father Adam, to bring back all that can be saved and placed where they can receive salvation and a glory in some kingdom. This Priesthood has to do it; this Priesthood is for this purpose.
		\end{compactdesc}
	
% ----------------------
\subsection{Commentary and Studies} % *COMMENTARY *STUDIES
% ----------------------
	\label{PRIESTHOOD-COMMENTARY}

	% % -----
	% \subsubsection{Key Points}
	% % -----
		% \label{PRIESTHOOD-KEYS}
	
		% \begin{compactitem}
			% \item
		% \end{compactitem}
		
	% -----
	\subsubsection{Definitions of the priesthood, or the proper way to understand what the priesthood is}
	% -----
		\label{PRIESTHOOD-definition}
		
		See \hyperref[ANOMIA]{\grl{ἀνομία}}; \hyperref[1CORINTHIANS-15-56]{1 Corinthians 15:56}; \hyperref[DC29-36]{DC 29:36} and notes there; \hyperref[MOSES-4-1]{Moses 4:1--3} and notes there; 
		
		Probably less important, but see \hyperref[DC53-3]{DC 53:3} and the older versions of it (the ordinances were associated with being an elder, not the ordination); \hyperref[DC68-1]{DC 68:1--3}; 
		
	% -----
	\subsubsection{Powers of the priesthood}
	% -----
		\label{PRIESTHOOD-powers}
	
		\begin{compactitem}
		
			\item The ``power on earth to forgive sins'' (\hyperref[MATTHEW-9-6]{Matthew 9:6})
			
			\item The ``power to make intercession for the children of men'' (\hyperref[MOSIAH-15-8]{Mosiah 15:8})
			
			\item The ``power of the resurrection'' (\hyperref[2NEPHI-10-25]{2 Nephi 10:25}); see also \hyperref[JACOB-4-11]{Jacob 4:11}, \hyperref[MORONI-7-41]{Moroni 7:41}, \hyperref[2NEPHI-9-6]{2 Nephi 9:6}, \hyperref[ECCLESIASTES-8-8]{Ecclesiastes 8:8}, \hyperref[PHILIPPIANS-3-10]{Philippians 3:10} (\gr{δύναμις}),
			
			\item The ``power of the atonement'' (\hyperref[2NEPHI-10-25]{2 Nephi 10:25})
			
			\item ``Power over death'' (\hyperref[DC7-2]{DC 7:2}; note, though, that this ``power over death'' part is not in the original manuscript version nor in the BC. See also \hyperref[REVELATIONNT-1-18]{Revelation 1:18})
			
			\item The ``power to speak'' (\hyperref[DC11-10]{DC 11:10--11})
			
			\item The ``power to obtain eternal life'' (\hyperref[DC45-8]{DC 45:8})
			
			\item The ``power to become the sons of God'' (\hyperref[DC11-30]{DC 11:30})
				\begin{compactitem}
					\item See also \hyperref[NEWBIRTH-sons]{this study} for more
				\end{compactitem} 
			
			\item The ``power to overcome all things which are not ordained of him'' (\hyperref[DC50-35]{DC 50:35})
			
			\item The ``power to command the waters'' (\hyperref[DC61-27]{DC 61:27}; see also \hyperref[JACOB-4-6]{Jacob 4:6} and \hyperref[MARK-4-39]{Mark 4:39})
			
			\item The ``power to seal up unto eternal life (\hyperref[DC68-12]{DC 68:12})
			
		\end{compactitem}
		
	% -----
	\subsubsection{The difference between the Aaronic and Melchizedek priesthoods}
	% -----
		\label{PRIESTHOOD-difference}
		
		See \hyperref[DC13-1]{DC 13}, along with references there to other passages describing the Aaronic priesthood;
		
	% -----
	\subsubsection{The ``order of Enoch''}
	% -----
		\label{PRIESTHOOD-enoch}
		
		See \hyperref[DC76-57]{DC 76:57}, and also the \hyperref[JST-OTR1-GENESIS-14]{JST for Genesis 14}.
		
	% -----
	\subsubsection{``Kingdom of God'' as a way of talking about the priesthood}
	% -----
		% \label{PRIESTHOOD-kingdom}
		
		Perhaps understand the term ``kingdom of God'' as being about a certain system of laws that becomes the way that people are governed or govern themselves; see \hyperref[LUKE-16-16]{Luke 16:16} and \hyperref[DC88-36]{DC 88:36}, among other verses.
		
	% -----
	\subsubsection{The in-dwelling relationship of the priesthood and new birth}
	% -----
		
		See this idea at this study of becoming the ``\hyperref[NEWBIRTH-sons]{sons of God}'' and other material around there; see also the ``\hyperref[NEWBIRTH-children-christ]{children of Christ}'' idea as well, which is very important for the priesthood and the name of the priesthood in 107:3.